

So the inverse FFT algorithm is similar to FFT, except that we use \( \omega_n^{-1} \) instead of \( \omega_n \), and divide the final result by \( n \).

\begin{algorithm}[H]
	\KwIn{$(y_0, y_1, \ldots, y_{n-1})$ where \( Y = (y_0, y_1, \ldots, y_{n-1}) \) is the point-value representation of \( A(x) \) at \( n \)-th roots of unity.}
	\KwOut{Coefficients representation of \( A(x) \).}
	$A \gets \text{FFT}(Y, \omega_n^{-1}) / n$\;
	\Return{\( A \)}\;
	\caption{Inverse FFT}
\end{algorithm}

\begin{theorem}
	Both FFT and Inverse FFT run in \( O(n\log n) \) time.
\end{theorem}
The proof is similar to the previous ones, so we omit it here.

\begin{note}
	In lecture, we spend lots of time showing how $3\times 3 = 9$ using FFT, and it takes an entire whiteboard.

	We don't cover it here as one can easily construct the example by themselves using the algorithms above.
\end{note}

\section{Examples using Convolution}

\begin{question}
	Given a set $A$ of $n$ positive integers with $\max(A) = M$, find $a, b, c\in A$ such that $a + b = c$.
\end{question}

A common solution is iterating through all pairs \( (a,b) \) and checking if \( a+b \in A \), which takes \( O(n^2) \) time.

Here we present a solution using convolution and FFT that runs in \( O(M \log M) \) time:

Let $V(i) = 1$ if $i \in A$, and $0$ otherwise. Set $\mu = V * V$. Then $\mu(k) = \sum_{i\leq k} V(i)V(k-i)$ counts the number of pairs $(a,b) \in A \times A$ such that $a + b = k$. Thus, if there exists $c \in A$ such that $\mu(c) > 0$, then there exist $a,b \in A$ such that $a + b = c$, and performing a linear scan will find such a triple $(a,b,c)$.

\begin{question}
	In set $A\subseteq [W]$ find solution for $x+2y-3z = 0$.
\end{question}

The solution will be similar to the previous question, the readers are encouraged to think about it by themselves.
